\documentclass[12pt]{article}
\usepackage[a4paper,margin=1in]{geometry}
\usepackage{amsmath}
\usepackage{amssymb}
\usepackage{siunitx}
\usepackage{enumitem}
\usepackage{multicol}

\usepackage[T1]{fontenc} % higher quality font encoding
\usepackage[sfdefault]{courier} % load the font and set it to default
% \renewcommand*\familydefault{\sfdefault} %sansserify


\begin{document}
\pagenumbering{gobble}

\title{Conversions: Decimals, Percentages, and Fractions}
\date{}
\maketitle

\begin{enumerate}[label=\textbf{Q\arabic*.}]


\item Write the following decimals as a fraction in simplest form and as a percentage:
\begin{multicols}{2}
\begin{enumerate}[label=(\alph*)]
    \item \( 0.125 \)
    \item \( 0.203125 \)
\end{enumerate}
\end{multicols}


\item Convert the following fractions into decimals and percentages:
\begin{multicols}{2}
\begin{enumerate}[label=(\alph*)]
    \item \( \frac{7}{20} \)
    \item \( \frac{33}{320} \)
\end{enumerate}
\end{multicols}

\item Express the following percentages as decimals and fractions in simplest form:
\begin{multicols}{2}
\begin{enumerate}[label=(\alph*)]
    \item \( 17.5\% \)
    \item \( 57.8125\% \)
\end{enumerate}
\end{multicols}

\end{enumerate}


\bigskip


\begin{enumerate}[label=\textbf{Q\arabic*.}, start=4]

\item

\textbf{Basis points} are units used in finance to describe small changes of percentages. A basis point is a percent of a percent:  
\[
1 \text{ basis point} = 0.01\%
\]

They are often used to describe changes in interest rates or yields.

For example, an increase of \( 25 \) basis points means the interest rate goes up by \( 0.25\% \).

\medskip





\begin{enumerate}[label=(\alph*)]
    \item Write one basis point as a decimal and a fraction in simplest form.
    \item Write 175 basis points as a decimal and a fraction in simplest form.
\end{enumerate}

\bigskip
    \textbf{Gilts} are UK government bonds. When you buy a gilt, you are lending money to the government.  

    The UK government currently issues 30-year gilts offering a simple return (not compounding) of $5.355\%$ per year. 

    They use these to borrow money to fund investment. Each year they pay a \textbf{coupon} equal to $5.355\%$ of the amount borrowed. At the end of the 30 years they pay back the original amount borrowed. 
    
    Other durations (1-year, 5-year, 10-year) of Gilts are also sold.

\medskip

\begin{enumerate}[label=(\alph*), resume]
    \item Find $5.355\%$ as a decimal and a fraction in its simplest form.
    \item Calculate the total cost for the government to borrow £1000 over 30 years, including payback of the original amount.
    \item Suppose the coupon payment of a 30 year gilt rises \( 43.5 \) basis points. How much more does it cost the government to borrow £1000 over 30 years?
    % \item Gilts offer only simple interest, how might an investor achieve compounding returns, while still only investing in Gilts?
\end{enumerate}

\end{enumerate}

\end{document}
